\begin{titlepage}

\begin{center}

\begin{figure}
\begin{minipage}[H]{6cm}
\centering
\includegraphics[width=0.8\linewidth]{img/university_logo}
\end{minipage}
\hfill
\begin{minipage}[H]{6cm}
\centering
% Company/Project logo placeholder
\textbf{\Large LogChirpy}
\end{minipage}
\end{figure}

\vspace*{0.8cm}

\textbf{Hochschule Reutlingen Fakultät Informatik}\\
\textbf{Mobile Computing Project"} \\
\vspace*{0.2cm}
{\large \textbf{SoSe 2025\\}}
\vspace*{0.4cm}
Prof. Dr. Martinez\\
\vspace*{0.8cm}
{\large \textbf{- Final Report -}}

\textbf{{\large {\LARGE LogChirpy - Ornithological Archival App}}}\\
\textbf{{\large Bird Watching Application with AI-Powered Recognition}}\\

\noindent\rule{\textwidth}{2pt}
\vspace*{0.6cm}

\textbf{vorgelegt von:\\}
{\large \textbf{Martin Lauterbach\\}}
{\large \textbf{Luis Wehrberger\\}}
{\large \textbf{Youmna Samouneh\\}}

\vspace*{0.6cm}

Contact addresses:\\
martin.lauterbach@student.reutlingen-university.de\\
luis.wehrberger@student.reutlingen-university.de\\
youmna.samouneh@student.reutlingen-university.de\\

\vspace*{0.6cm}

Submitted on: \today

\vspace*{1cm}

\textbf{Abstract}\\
\vspace*{0.3cm}
\begin{minipage}{0.8\textwidth}
\small
This research presents LogChirpy, a comprehensive mobile ornithological documentation application developed using React Native and Expo frameworks. The application implements advanced artificial intelligence methodologies for automated avian species identification through integrated visual and acoustic analysis systems. The architecture employs local machine learning models for real-time species recognition, coupled with a sophisticated archival framework that supports both offline operation and optional cloud synchronization capabilities. The system incorporates multilingual functionality across six languages and maintains a comprehensive taxonomic database encompassing over 30,000 global avian species. LogChirpy addresses the requirements of both amateur naturalists and professional ornithologists by providing an intuitive platform for systematic documentation and collaborative sharing of ornithological observations.
\end{minipage}

\end{center}

\end{titlepage}